% filename: chap-conclusion-FINAL-v2.tex
\chapter{Conclusion and Future Work}
\label{chap:conclusion}
\setlength{\emergencystretch}{1.5em}   % prevents right-margin overflow in this chapter

% ============================================================
\section{Summary of the Work}
\label{sec:conclusion_summary}
% ============================================================

This thesis investigated whether a unified reinforcement learning agent,
jointly optimising energy dispatch and diesel inventory ordering, can
improve the operational reliability of off-grid and weak-grid telecom
base stations relative to decoupled classical approaches.  The work was
motivated by a practical gap in the literature: while prior work on
RL-based energy management of remote base stations routinely models
battery dispatch and diesel generator scheduling, it consistently omits
the upstream supply chain decision --- when and how much diesel to order
--- despite this being a primary cause of service interruption in
fuel-constrained deployments.

The study framed this gap as a constrained Markov decision process
(CMDP) in which a single policy controls diesel generator on/off state
and periodic diesel order quantity.  Battery dispatch is determined by
an internal environment rule rather than a direct RL action.  The
inventory pipeline --- litres of diesel in transit under stochastic
lead times --- was incorporated as an explicit state variable, enabling
the agent to anticipate impending stockouts and adjust both dispatch
conservatism and order timing accordingly.  This formulation, denoted
\emph{RLInv}, was trained using Proximal Policy Optimisation (PPO)
with action masking to enforce hard feasibility constraints, together
with a fixed penalty term on soft constraint violations during training.

The thesis also implemented a hierarchical decomposition baseline
(TrackB) in which diesel ordering is governed by a classical
$(s,S)$ inventory policy and only dispatch is learned by RL.  This
directly addressed the timescale decomposition critique raised in the
Phase-I review: the empirical comparison between RLInv and TrackB
constitutes a direct test of whether joint optimisation provides
measurable benefit over principled decomposition.

All experiments used the ITU~2024 Smart Energy Scheduling dataset
(ten remote base station sites, 60~days each, hourly resolution),
with a synthetic lead-time overlay applied only to the in-transit
diesel state variable~$\mathrm{Pipe}_t$, which is not present in
the original dataset.  Three pre-registered hypotheses were tested
across a structured ablation study of eight policy variants.

% ============================================================
\section{Key Results and Contributions}
\label{sec:conclusion_results}
% ============================================================

\subsection{Principal Empirical Findings}

The central empirical finding is that RLInv achieves a mean
Expected Energy Not Served (EENS) of $488.5 \pm 92.4$\,kWh per
episode on constrained sites under normal logistics, compared to
$730.9 \pm 14.0$\,kWh for both the hierarchical RL baseline (TrackB)
and the fixed-ordering ablation (A6).  Relative to RLInv, comparator
EENS for both baselines is 49.6\% higher.  This difference is
statistically significant ($p=0.042$, Cohen's $d=3.67$, Welch's
$t$-test, $n=3$ seeds).

Three hypotheses were evaluated:

\sloppy
\begin{description}
  \item[H1 (Joint ordering).] \emph{Supported} on constrained sites
  under normal logistics.  The principal performance difference arises
  from the ordering mechanism rather than the dispatch architecture:
  A6 and TrackB produce numerically identical EENS ($730.9$\,kWh,
  $p=1.0$), confirming that dispatch policy architecture is not the
  determining factor once ordering is fixed.

  \item[H2 (Action masking).] \emph{Partially supported.}  Removing
  masking (A7) produced a between-seed standard deviation of
  $358.9$\,kWh versus $92.4$\,kWh for RLInv, consistent with
  substantially less stable outcomes across training runs.  No hard
  feasibility violations were recorded at evaluation for either
  policy, suggesting masking's primary benefit is training
  consistency rather than deployment safety enforcement.  The mean
  EENS difference was not statistically significant ($p=0.525$,
  $d=0.61$) at $n=3$ seeds.

  \item[H3 (Inventory observation).] \emph{Directionally supported}
  under normal logistics: A5 exhibits 15.6\% higher EENS than RLInv
  ($p=0.429$, $d=0.72$), but the result is \emph{not statistically
  conclusive} at $n=3$ seeds.  Under delayed logistics, the benefit
  disappears and slightly reverses ($-2.2\%$, $p=0.398$), indicating
  that the effect is distribution-specific.
\end{description}
\fussy

The most important secondary finding concerns robustness to logistics
stress.  Under delayed delivery ($\mu=48$\,h lead time), RLInv's EENS
increases by $+596$\,kWh ($+122\%$) --- the largest degradation of any
policy evaluated.  By contrast, A6 and TrackB degrade by only
$+323$\,kWh ($+44\%$), while B0 degrades by a similar amount.  The
$(s,S)$ safety stock therefore provides structural robustness to
lead-time uncertainty that the RL agent does not learn from
single-scenario training.  This finding is the primary motivation
for multi-scenario training in Phase~III.

Cross-site evaluation confirmed that RLInv does not degrade
meaningfully relative to the best rule-based baseline on easy or
medium sites.  Reliability improvements concentrate on the two
hardest, fuel-constrained sites (site~5 and site~10), consistent with
the expectation that inventory management matters most where energy
resources are scarce.

\subsection{Novel Contributions}

This thesis makes the following contributions to the literature on
RL-based energy management of remote telecom infrastructure:

\begin{enumerate}

  \item \textbf{CMDP formulation with explicit inventory pipeline state.}
  The inclusion of in-transit diesel ($\mathrm{Pipe}_t$) as a state
  variable, alongside fuel level, battery SoC, and grid availability,
  enables anticipatory ordering decisions that react to lead-time
  uncertainty before a stockout occurs.  To the best of the author's
  knowledge, this state augmentation has not been applied to telecom
  base station energy management in prior work.

  \item \textbf{Empirical test of the coupling hypothesis.}
  Rather than assuming that joint optimisation is superior to
  decomposition, this work constructs a controlled ablation
  (RLInv~vs.~A6~vs.~TrackB) that directly measures the value of
  coupling.  The finding that joint ordering is beneficial under
  normal logistics, but that decomposition is \emph{more robust}
  under logistics stress, is a nuanced result with practical
  implications for deployment strategy.

  \item \textbf{Action masking for training stability in hybrid action
  spaces.}  The use of MaskablePPO with analytically derived action
  bounds provides a computationally lightweight mechanism for
  maintaining feasibility during training in a space combining
  discrete generator commitment and ordering decisions.  The ablation
  demonstrates that masking primarily reduces outcome variance rather
  than enforcing hard constraints at deployment.

  \item \textbf{Reproducible open benchmark.}
  The ITU/Zindi dataset, simulator, and evaluation pipeline are
  publicly available, enabling direct comparison with Ma \&
  Pan~(2025) and future work.  The structured ablation design
  (eight policy variants, three seeds, two logistics scenarios,
  ten sites) provides a reusable evaluation framework for
  inventory-aware energy management research.

\end{enumerate}

% ============================================================
\section{Critique and Limitations}
\label{sec:conclusion_limitations}
% ============================================================

A number of constraints imposed by the 14-day execution window and
the CPU-only compute environment limit the generalisability of the
present results.  These are documented as a precise statement of scope
rather than as weaknesses to be apologised for; each directly motivates
the Phase~III plan in Section~\ref{sec:conclusion_future}.

\sloppy
\begin{description}

  \item[Statistical power.] All hypothesis tests are conducted at
  $n=3$ seeds.  H2 and H3 are directionally consistent with the
  proposed mechanisms but do not reach statistical significance.
  Definitive conclusions on the contribution of action masking and
  inventory observation require more seeds, which was not achievable
  within the time budget.

  \item[Synthetic lead-time distribution.] The in-transit state
  variable $\mathrm{Pipe}_t$ is generated from a synthetic
  lognormal distribution ($\mu=24$\,h, $\sigma=0.6$) because the
  ITU dataset does not record diesel delivery event timestamps.
  Real logistics patterns --- including seasonal supply disruptions
  and road access constraints --- are not represented.

  \item[Single logistics scenario during training.] RLInv is
  trained exclusively under normal logistics and evaluated on
  both normal and delayed scenarios.  The $+122\%$ EENS degradation
  under delay indicates that the learned ordering strategy is
  optimised for the training distribution and does not generalise
  to out-of-distribution lead-time stress without multi-scenario
  training.

  \item[Reduced training scope.] To meet the compute budget,
  training episodes were shortened to 30~days, the discount factor
  was reduced to $\gamma=0.995$ (effective horizon $\approx
  200$\,h), and order quantity was discretised to five levels.
  These simplifications limit the policy's ability to learn
  long-horizon inventory strategies.

  \item[Ablation coverage gaps.] A7 (no action masking) was not
  evaluated under the delayed logistics scenario, and the A2--A4
  ablations (individual safety component removal) were not executed
  within the available time.  The H2 finding therefore rests on a
  limited ablation coverage.

  \item[RL-only algorithmic comparison.] The present results
  establish RL-vs-RL comparisons across the ablation study but do
  not yet include a value-based RL ordering baseline, which would
  clarify whether gains arise from the inventory-aware formulation
  or the specific policy-gradient learning mechanism.

\end{description}
\fussy

% ============================================================
\section{Future Work and Phase-III Plan}
\label{sec:conclusion_future}
% ============================================================

The results in Chapter~\ref{chap:results} show that RLInv improves
reliability on constrained sites under the nominal logistics regime,
but its performance degrades substantially under delayed replenishment.
Accordingly, the next phase of this research will focus on improving
robustness to logistics uncertainty, strengthening statistical
confidence, and expanding algorithmic comparison within the existing
simulator framework.  The four extensions below are scoped to a
realistic continuation of four to six weeks.

\subsection*{1.\ Multi-scenario training for logistics robustness}

The most immediate extension is to train RLInv across multiple
replenishment regimes rather than a single nominal scenario.
Concretely, the lead-time distribution $\mathrm{LogNormal}(\mu,{}\allowbreak\sigma)$
can be varied across episodes by sampling $\mu$ uniformly from the
range $[12, 72]$\,h during training, while evaluation is performed on
fixed held-out regimes including the delayed scenario used in
Chapter~\ref{chap:results}.  This domain-randomisation approach requires
no structural changes to the simulator and directly tests whether
exposure to logistics uncertainty during training can close the
$+122\%$ EENS degradation gap observed under delayed delivery.

A related modification is to introduce modest stochasticity in
individual delivery realisations --- partial delays and occasional
split deliveries --- while keeping the existing environment interface
unchanged.  This provides a more realistic stress-test of the ordering
policy without requiring a redesign of the CMDP formulation.

\subsection*{2.\ Expanded statistical evaluation}

The present ablation results rest on three training seeds, which is
sufficient to report directional effects but insufficient to confirm
statistical significance for H2 and H3.  A near-term extension is
therefore to increase the seed count to ten or more and rerun the
constrained-site evaluation, reporting 95\% confidence intervals
alongside point estimates.  This would allow the H2 (action masking)
and H3 (inventory observation) verdicts to be revisited with greater
confidence, or to be definitively classified as effects that are too
small to detect reliably at the current training scale.

The evaluation protocol can further be strengthened by testing policy
behaviour under additional stress scenarios, including extended grid
outage stretches and more variable replenishment delays, to probe
rare but operationally critical conditions.

\subsection*{3.\ Value-based algorithmic baseline for ordering}

The present thesis focuses on a policy-gradient approach (PPO) for
both RL action components.  A practical next step is to add one
value-based RL baseline specifically for the diesel ordering decision
--- for example, Double DQN over the five-level discrete order action
space --- within the existing simulator.  This comparison would clarify
whether the reliability gains demonstrated by RLInv arise from the
inventory-aware state formulation itself or also depend on the
policy-gradient learning mechanism.  The implementation cost is low
because the simulator interface and evaluation pipeline require no
modification.

\subsection*{4.\ Cross-site and cross-distribution generalisation}

The current evaluation treats each site as an independent training
target.  A final extension is to test whether the learned policy
captures general operational structure or depends heavily on
site-specific patterns.  This can be assessed by training on a subset
of ITU sites and evaluating on held-out sites, or by training under
one logistics regime and evaluating under another.  Both experiments
reuse the existing dataset and simulator, and would strengthen
the practical claim that inventory-aware RL can generalise beyond
the exact conditions represented during training.

\subsection*{Indicative Phase-III timeline}

\begin{itemize}
  \item \textbf{Weeks 1--2:} implement multi-scenario training across
  normal, delayed, and mixed logistics regimes; regenerate the
  constrained-site evaluation and compare against current Phase-II
  results.

  \item \textbf{Weeks 3--4:} increase seed count to ten and rerun the
  full ablation study; revise H2 and H3 verdicts using confidence
  intervals and expanded seed-level statistics.

  \item \textbf{Weeks 5--6:} add the value-based ordering baseline and
  conduct cross-site / cross-distribution generalisation experiments;
  consolidate all Phase-III results into a revised evaluation chapter.
\end{itemize}

\subsection*{Longer-term research directions}

Beyond the six-week Phase-III scope, several directions follow
naturally from the problem formulation and are worth identifying
for the record, even where they are not immediate deliverables.

\textbf{Classical optimisation comparison.}
The present results establish RL-vs-RL comparisons across the ablation
study, but a complete evaluation requires comparison against
deterministic baselines used in commercial deployments.  Model
Predictive Control (MPC) with a rolling forecast horizon is the
natural comparator: it handles the dispatch problem well but typically
relies on fixed safety-stock rules for the ordering decision, exactly
the gap that RLInv targets.  Demonstrating that inventory-aware RL
matches or exceeds MPC under stochastic lead times --- where
certainty-equivalent planning is most vulnerable to lead-time
uncertainty --- would be the strongest possible validation of the
coupling hypothesis.

\textbf{Carbon and cost multi-objective reporting.}
The objective function defined in Chapter~\ref{chap:problem} includes
a carbon cost term $c_{\mathrm{CO}_2}$, but the evaluation in
Chapter~\ref{chap:results} reports only EENS and diesel consumption.
A natural extension is to report Pareto frontiers in the
EENS--CO$_2$ plane, testing whether the carbon penalty can reduce
emissions without increasing service interruptions.  This requires no
simulator change --- only an extension of the evaluation pipeline ---
and would directly address sustainability framing that is increasingly
required by telecom regulators in low- and middle-income countries.

\textbf{Solar forecast integration.}
RLInv currently observes the realised PV output at each timestep.
In practice, day-ahead solar irradiance forecasts are available from
numerical weather prediction services.  Augmenting the state space
with a short-horizon forecast vector would allow the agent to
anticipate low-PV periods and pre-position battery and diesel
inventory accordingly.  This is a well-established improvement
direction in RL-based energy management and would be straightforward
to implement within the existing Gymnasium environment.

\textbf{Multi-tower shared-supplier formulation.}
In realistic deployments, multiple base stations share a common
diesel supplier and compete for delivery slots, particularly during
natural disasters or supply disruptions.  The single-tower CMDP
formulated here cannot capture this coordination problem.  A
multi-agent extension, in which each tower is an independent agent
but lead times are coupled through shared logistics capacity, would
address this gap and is the most direct route to operationally
relevant deployment at scale.

% ============================================================
\section{Closing Remarks}
\label{sec:conclusion_closing}
% ============================================================

This thesis set out to answer a specific question: does making the
inventory pipeline visible to the RL controller --- letting it see
how much diesel is in transit and decide when to order ---
measurably reduce service interruptions at fuel-constrained
telecom base stations?  Under the logistics conditions the agent
was trained on, the answer is yes, and the improvement is large
and statistically significant.  Under out-of-distribution logistics
stress, the advantage disappears, revealing a robustness gap that
is well-understood theoretically and correctable in Phase~III.

The thesis also demonstrates that the scientific question
matters as much as the engineering result.  The decision to frame
the study as a hypothesis test --- rather than as a claim that
RL is superior to classical methods --- allowed the negative and
null results (H2 and H3 inconclusive, H1 benefit disappears under
delay) to contribute meaningfully rather than being suppressed.
The finding that timescale decomposition is sufficient when logistics
are stressed is an honest and useful result for practitioners
deciding how to deploy energy management software at remote sites.

Energy security at remote telecom towers is not a niche problem.
In many low- and middle-income countries, base station uptime
determines access to emergency services, financial transfers, and
public health communication.  A fuel stockout at the wrong site
at the wrong time has consequences far beyond the cost of a
missed call.  Building controllers that are both reliable under
normal conditions and robust to supply disruptions is, therefore,
not only a technical challenge but also an infrastructural necessity.